\documentclass[../../../include/open-logic-section]{subfiles}

\begin{document}

\olfileid[pt]{sth}{ordinals}{ordtype}
\olsection{Ordinais como Tipos de Ordem}

Armados com a Substituição, e portanto trabalhando agora em $\ZFminus$,
podemos finalmente provar o resultado a que vimos aspirando:

\begin{thm}\ollabel{thmOrdinalRepresentation}
Toda a boa ordenação é isomorfa a um único ordinal.
\end{thm}

\begin{proof}
Seja $\tuple{A, <}$ uma boa ordenação. Por
\olref[basic]{ordisoidentity}, é isomorfa no máximo a um
ordinal. Logo, por redução ao absurdo, suponha que $\tuple{A, <}$ não é
isomorfa a \emph{nenhum} ordinal. Começaremos por ``tornar $\tuple{A, <}$ o mais pequena possível''. Em detalhe: se algum segmento inicial próprio $\tuple{A_a,
<_a}$ não for isomorfo a nenhum ordinal, existe um $a \in A$ mínimo com
essa propriedade; então seja $B = A_a$ e $\mathord{\lessdot} =
\mathord{<_a}$. Caso contrário, seja $B = A$ e $\mathord{\lessdot} =
\mathord{<}$.

Por definição, todo o segmento inicial próprio de $B$ é isomorfo a
algum ordinal, que é único como acima. Logo,
pela Substituição, o conjunto seguinte existe, e é uma função:
\[
	f = \Setabs{\tuple{\beta, b}}{b \in B\text{ e }\ordeq{\beta}{\tuple{B_b, \lessdot_b}}}
\]
Para completar a redução ao absurdo, mostraremos que $f$ é um isomorfismo
$\alpha \to B$, para algum ordinal $\alpha$.

É óbvio que $\ran{f}
= B$. E por \olref[iso]{lemordsegments}, $f$ preserva a ordenação,
isto é, $\gamma \in \beta$ sse $f(\gamma) \lessdot f(\beta)$. Para mostrar que $\dom{f}$ é um ordinal, por
\olref[basic]{corordtransitiveord} basta mostrar que
$\dom{f}$ é transitivo. Fixe então $\beta \in \dom{f}$, isto é,
$\ordeq{\beta}{\tuple{B_b, \lessdot_b}}$ para algum $b$. Se $\gamma \in
\beta$, então $\gamma \in \dom{f}$ por
\olref[iso]{wellordinitialsegment}; generalizando, $\beta \subseteq
\dom{f}$.
\end{proof}

Este resultado autoriza a definição seguinte, que queríamos apresentar
desde \olref[vn]{sec}:

\begin{defn}
Se $\tuple{A, <}$ é uma boa ordenação, então o seu tipo de ordem,
$\ordtype{A, <}$, é o único ordinal $\alpha$ tal que
$\ordeq{\tuple{A, < }}{\alpha}$.
\end{defn}

Além disso, esta definição autoriza dois princípios convenientes:

\begin{cor}\ollabel{ordtypesworklikeyouwant}
Sendo $\tuple{A, <}$ e $\tuple{B, \lessdot}$ boas ordenações: 
\begin{align*}
	\ordtype{A, <} = \ordtype{B, \lessdot}&\text{ sse }\ordeq{\tuple{A, <}}{\tuple{B, \lessdot}}\\
	\ordtype{A, <} \in \ordtype{B, \lessdot}&\text{ sse }\ordeq{\tuple{A, <}}{\tuple{B_b, \lessdot_b}}\text{ para algum }b \in B
\end{align*}
\end{cor}

\begin{proof}
A identidade vale por \olref[basic]{ordisoidentity}. Para provar
a segunda afirmação, seja $\ordtype{A, <} = \alpha$ e $\ordtype{B,
\lessdot} = \beta$, e seja $f \colon \beta \to \tuple {B, \lessdot}$
o nosso isomorfismo. Então:
\begin{align*}
	\alpha \in \beta&\text{ sse }\funrestrictionto{f}{\alpha} \colon \alpha \to B_{f(\alpha)}\text{ é um isomorfismo}\\
	&\text{ sse }\ordeq{\tuple{A, <}}{\tuple{B_{f(\alpha)}, \lessdot_{f(\alpha)}}}\\
	&\text{ sse }\ordeq{\tuple{A, <}}{\tuple{B_b, \lessdot_b}}\text{ para algum $b \in B$}
\end{align*}
por \olref[iso]{ordisounique},
\olref[iso]{wellordinitialsegment}, e
\olref[basic]{ordissetofsmallerord}.
\end{proof}

\end{document}
